\documentclass[10pt,a4paper]{amsart}
\usepackage{amsmath,amssymb,amsthm,graphicx}
% Retained minimal preamble for the Stacks Project chapter varieties.tex.
% Verbatim excerpt of the chapter preamble: Xy-pic setup (native preamble lines 14-18), the native theorem
% declarations with the mathematical macros (native preamble lines 154-217), and the project own no-op
% definitions of the reference/slogan/history display environments (native wrapper.tex lines 1-3) that the
% counted statement uses for its slogan. Copied mechanically; nothing added, removed or redefined.
% For commutative diagrams we use Xy-pic
\usepackage[all]{xy}

% We use 2cell for 2-commutative diagrams.
\xyoption{2cell}
% Theorem environments.
%
\theoremstyle{plain}
\newtheorem{theorem}[subsection]{Theorem}
\newtheorem{proposition}[subsection]{Proposition}
\newtheorem{lemma}[subsection]{Lemma}

\theoremstyle{definition}
\newtheorem{definition}[subsection]{Definition}
\newtheorem{example}[subsection]{Example}
\newtheorem{exercise}[subsection]{Exercise}
\newtheorem{situation}[subsection]{Situation}

\theoremstyle{remark}
\newtheorem{remark}[subsection]{Remark}
\newtheorem{remarks}[subsection]{Remarks}

\numberwithin{equation}{subsection}

% Macros
%
\def\lim{\mathop{\mathrm{lim}}\nolimits}
\def\colim{\mathop{\mathrm{colim}}\nolimits}
\def\Spec{\mathop{\mathrm{Spec}}}
\def\Hom{\mathop{\mathrm{Hom}}\nolimits}
\def\Ext{\mathop{\mathrm{Ext}}\nolimits}
\def\SheafHom{\mathop{\mathcal{H}\!\mathit{om}}\nolimits}
\def\SheafExt{\mathop{\mathcal{E}\!\mathit{xt}}\nolimits}
\def\Sch{\mathit{Sch}}
\def\Mor{\mathop{\mathrm{Mor}}\nolimits}
\def\Ob{\mathop{\mathrm{Ob}}\nolimits}
\def\Sh{\mathop{\mathit{Sh}}\nolimits}
\def\NL{\mathop{N\!L}\nolimits}
\def\CH{\mathop{\mathrm{CH}}\nolimits}
\def\proetale{{pro\text{-}\acute{e}tale}}
\def\etale{{\acute{e}tale}}
\def\QCoh{\mathit{QCoh}}
\def\Ker{\mathop{\mathrm{Ker}}}
\def\Im{\mathop{\mathrm{Im}}}
\def\Coker{\mathop{\mathrm{Coker}}}
\def\Coim{\mathop{\mathrm{Coim}}}

% Boxtimes
%
\DeclareMathSymbol{\boxtimes}{\mathbin}{AMSa}{"02}

%
% Macros for moduli stacks/spaces
%
\def\QCohstack{\mathcal{QC}\!\mathit{oh}}
\def\Cohstack{\mathcal{C}\!\mathit{oh}}
\def\Spacesstack{\mathcal{S}\!\mathit{paces}}
\def\Quotfunctor{\mathrm{Quot}}
\def\Hilbfunctor{\mathrm{Hilb}}
\def\Curvesstack{\mathcal{C}\!\mathit{urves}}
\def\Polarizedstack{\mathcal{P}\!\mathit{olarized}}
\def\Complexesstack{\mathcal{C}\!\mathit{omplexes}}
% \Pic is the operator that assigns to X its picard group, usage \Pic(X)
% \Picardstack_{X/B} denotes the Picard stack of X over B
% \Picardfunctor_{X/B} denotes the Picard functor of X over B
\def\Pic{\mathop{\mathrm{Pic}}\nolimits}
\def\Picardstack{\mathcal{P}\!\mathit{ic}}
\def\Picardfunctor{\mathrm{Pic}}
\def\Deformationcategory{\mathcal{D}\!\mathit{ef}}
\newenvironment{reference}{}{}
\newenvironment{slogan}{}{}
\newenvironment{history}{}{}
\usepackage{cleveref}
\crefname{originallemma}{original lemma}{original lemmas}
\Crefname{originallemma}{Original lemma}{Original lemmas}
\makeatletter
\newlabel{lemma-surjective-pic-birational-finite}{{0C0T}{source}{}{sourcepointer-1}{}}
\newlabel{lemma-surjective-pic-birational-finite@cref}{{[originallemma][0][]0C0T}{[1][0][]source}{}{}{}}
\makeatother
\title{Glueing invertible sheaves and compatible nonvanishing sections from closed subschemes meeting in a discrete set}
\author{The Stacks Project Authors}
\date{stacks-project:varieties.tex:local-snapshot-20261011; GNU FDL 1.2 or later}
\begin{document}
\maketitle
\noindent\textit{Editorial note.} All delivered mathematical text is the author's own, copied verbatim from the pinned Stacks Project chapter file varieties.tex; no statement or proof step was rewritten, reordered or removed.\par\medskip
\noindent\textit{Editorial note.} The counted span cites one cross-chapter label, recorded in assembly.external\_references with its canonical tag; the author explicitly says a direct proof is also given, and that direct proof is what the counted span contains.\par\medskip
\noindent\textit{Editorial note.} The proof cites no bibliography item inside the counted span.\par\medskip
\noindent\textit{Allowed standard prerequisites (editorial).}\begin{itemize}
\item Surjectivity of the Picard map for a finite birational morphism. Stacks Project, Varieties, Lemma lemma-surjective-pic-birational-finite (tag 0C0T). Cited as an alternative route to the existence of the glued invertible sheaf; the proof also gives a self-contained direct construction, which is the part counted here.
\end{itemize}\medskip
\ifdefined\hypertarget\hypertarget{sourcepointer-1}{}\fi
\noindent\textit{Source reference.} Original lemma 0C0T: sources/stacks-chow/tags.txt, source lines 15098--15098. Cross-chapter standard prerequisite cited by the author inside the counted span; it is used as a pointer and its own proof is not part of this item.\par\medskip
\begin{lemma}
\label{lemma-glue-invertible-sheaves}
Let $X$ be a scheme. Let $Z_1, \ldots, Z_n \subset X$ be closed
subschemes. Let $\mathcal{L}_i$ be an invertible sheaf on $Z_i$.
Assume that
\begin{enumerate}
\item $X$ is reduced,
\item $X = \bigcup Z_i$ set theoretically, and
\item $Z_i \cap Z_j$ is a discrete topological space for $i \not = j$.
\end{enumerate}
Then there exists an invertible sheaf $\mathcal{L}$ on $X$ whose restriction
to $Z_i$ is $\mathcal{L}_i$. Moreover, if we are given sections
$s_i \in \Gamma(Z_i, \mathcal{L}_i)$ which are nonvanishing at the
points of $Z_i \cap Z_j$, then we can choose $\mathcal{L}$ such
that there exists a $s \in \Gamma(X, \mathcal{L})$ with
$s|_{Z_i} = s_i$ for all $i$.
\end{lemma}

\begin{proof}
The existence of $\mathcal{L}$ can be deduced from
Lemma \ref{lemma-surjective-pic-birational-finite}
but we will also give a direct proof and we will use
the direct proof to see the statement about sections is true.
Set $T = \bigcup_{i \not = j} Z_i \cap Z_j$. As $X$ is reduced we have
$$
X \setminus T = \bigcup (Z_i \setminus T)
$$
as schemes. Assumption (3) implies $T$ is a discrete subset of $X$.
Thus for each $t \in T$ we can find an open $U_t \subset X$
with $t \in U_t$ but $t' \not \in U_t$ for $t' \in T$, $t' \not = t$.
By shrinking $U_t$ if necessary, we may assume that there exist isomorphisms
$\varphi_{t, i} : \mathcal{L}_i|_{U_t \cap Z_i} \to
\mathcal{O}_{U_t \cap Z_i}$. Furthermore, for each $i$ choose an open covering
$$
Z_i \setminus T = \bigcup\nolimits_j U_{ij}
$$
such that there exist isomorphisms
$\varphi_{i, j} : \mathcal{L}_i|_{U_{ij}} \cong \mathcal{O}_{U_{ij}}$.
Observe that
$$
\mathcal{U} : X = \bigcup U_t \cup \bigcup U_{ij}
$$
is an open covering of $X$. We claim that we can use the isomorphisms
$\varphi_{t, i}$ and $\varphi_{i, j}$ to define a $2$-cocycle with values
in $\mathcal{O}_X^*$ for this covering that defines $\mathcal{L}$ as
in the statement of the lemma.

\medskip\noindent
Namely, if $i \not = i'$, then $U_{i, j} \cap U_{i', j'} = \emptyset$
and there is nothing to do. For $U_{i, j} \cap U_{i, j'}$ we have
$\mathcal{O}_X(U_{i, j} \cap U_{i, j'}) =
\mathcal{O}_{Z_i}(U_{i, j} \cap U_{i, j'})$ by the first remark of the proof.
Thus the transition function for $\mathcal{L}_i$ (more precisely
$\varphi_{i, j} \circ \varphi_{i, j'}^{-1}$) defines the value of our
cocycle on this intersection.
For $U_t \cap U_{i, j}$ we can do the same thing.
Finally, for $t \not = t'$ we have
$$
U_t \cap U_{t'} = \coprod (U_t \cap U_{t'}) \cap Z_i
$$
and moreover the intersection $U_t \cap U_{t'} \cap Z_i$ is contained
in $Z_i \setminus T$. Hence by the same reasoning as before we see that
$$
\mathcal{O}_X(U_t \cap U_{t'}) =
\prod \mathcal{O}_{Z_i}(U_t \cap U_{t'} \cap Z_i)
$$
and we can use the transition functions for $\mathcal{L}_i$ (more precisely
$\varphi_{t, i} \circ \varphi_{t', i}^{-1}$) to define the value of
our cocycle on $U_t \cap U_{t'}$. This finishes the proof of existence
of $\mathcal{L}$.

\medskip\noindent
Given sections $s_i$ as in the last assertion of the lemma, in the argument
above, we choose $U_t$ such that $s_i|_{U_t \cap Z_i}$ is nonvanishing and
we choose $\varphi_{t, i}$ such that $\varphi_{t, i}(s_i|_{U_t \cap Z_i}) = 1$.
Then using $1$ over $U_t$ and $\varphi_{i, j}(s_i|_{U_{i, j}})$ over
$U_{i, j}$ will define a section of $\mathcal{L}$ which restricts
to $s_i$ over $Z_i$.
\end{proof}

\end{document}
